% META: {"id":"TCOMPL-BAYES-EX-004","chapitre":"TCOMPL-INFERENCE-BAYESIENNE","type_objet":"exercice","capacites":["TCOMPL-BAYES-C5"],"capacites_codes":["C5"],"methodes":[],"parcours":2,"competences":["calculer","raisonner"],"duree_min":15,"mode_creation":"ex_nihilo","sources_inspiration":[],"fichier_tex":"chapitres/TCOMPL-INFERENCE-BAYESIENNE/exercices/TCOMPL-BAYES-EX-004.tex","corrige_tex":"chapitres/TCOMPL-INFERENCE-BAYESIENNE/corriges/TCOMPL-BAYES-CO-004.tex","status":"approved","origin":"generated","status_history":["generated","audited","accepted","approved"],"release_acceptance":"RELEASE_OWNER_BATCH_ACCEPTANCE","acceptance_closure_digest":"sha256:9b3ccf9a81c5520fbb7e03b7b2d3e7bf2057a02834b6d908bf3be5f49f3b553f","acceptance_receipt_digest":"sha256:4fad86f0282e0fa4e0419cca1ad6379ae7af5a280c04a4a90880f90a1c044242"}
% BEGIN-VERIFY
% from sympy import Rational, N
% Se, Sp = Rational(90, 100), Rational(95, 100)
% def VPP(p):
%     return Se*p / (Se*p + (1-Sp)*(1-p))
% assert abs(N(VPP(Rational(1,100))) - 0.1538) < 0.001
% assert abs(N(VPP(Rational(20,100))) - 0.8182) < 0.001
% assert VPP(Rational(20,100)) > VPP(Rational(1,100))
% END-VERIFY
\begin{exercice}{TCOMPL-BAYES-EX-004}{2}{15}
Un test a une sensibilité $Se=0{,}90$ et une spécificité $Sp=0{,}95$.
\begin{enumerate}
  \item Calculer la $VPP$ de ce test appliqué à une population générale de prévalence $p=0{,}01$.
  \item Calculer la $VPP$ de ce même test appliqué à une population à risque, de prévalence $p=0{,}20$.
  \item Comparer les deux résultats et expliquer, en une phrase, pourquoi ce même test est utilisé différemment selon la population ciblée.
\end{enumerate}
\end{exercice}
